\documentclass{article}
\usepackage[T1]{fontenc}
\usepackage[utf8]{inputenc}
\usepackage{amsmath,amssymb}
\usepackage{xcolor}
\usepackage{tikz}
\usepackage[american]{circuitikz}
\usepackage[version=4]{mhchem}
\usepackage{chemfig}
\usetikzlibrary{arrows.meta,shapes.geometric}
\begin{document}
% ---------------------------------------------------------------
% Generated by FlowTikZ (flowcharts and data charts to LaTeX)
% Every coordinate below is derived from the editor diagram model.
% Safe to edit by hand, but regenerating will overwrite changes.
% ---------------------------------------------------------------
% Add these to your preamble:
%   \usepackage{tikz}
%   \usepackage{xcolor}
%   \usetikzlibrary{arrows.meta,shapes.geometric}
%   \usepackage{amsmath}
%   \usepackage{amssymb}
%   \usepackage[american]{circuitikz}
%   \usepackage[version=4]{mhchem}
%   \usepackage{chemfig}

% ---- Colour definitions ----
\definecolor{fc12161C}{HTML}{12161C}
\definecolor{fc1F2933}{HTML}{1F2933}

% ---- Style definitions ----
\tikzset{
  fcnode/.style={
    draw, inner sep=0pt, outer sep=0pt, align=center,
    line join=round, line cap=butt
  },
  fcpath/.style={line join=round, line cap=butt},
  fcedge/.style={line join=round, line cap=butt},
  fclabel/.style={inner sep=1.4pt, outer sep=0pt, align=center},
}
\tikzset{fcformula/.style={inner sep=0pt, outer sep=0pt}}
\tikzset{fccircuit/.style={line join=round, line cap=butt,
  /tikz/circuitikz/bipoles/thickness=1, /tikz/circuitikz/tripoles/thickness=1}}

\begin{tikzpicture}[x=1cm, y=1cm]
  % ---- Nodes ----
  % n466 - Process "Plain flowchart box"
  \node[fcnode, shape=rectangle, minimum width=3.2cm, minimum height=1.4cm, draw=fc1F2933,
        line width=0.91pt, fill=white, font=\fontsize{7.97pt}{9.96pt}\selectfont\sffamily,
        text=fc12161C, align=center, text width=2.88cm]
        (n466) at (3.6,-2.7) {Plain flowchart box};
  % n467 - Decision "Math $x$ in text?"
  \node[fcnode, shape=diamond, minimum width=3cm, minimum height=2cm, draw=fc1F2933,
        line width=0.91pt, fill=white, font=\fontsize{7.97pt}{9.96pt}\selectfont\sffamily,
        text=fc12161C, align=center, text width=1.5cm]
        (n467) at (3.5,-5.8) {Math \$x\$ in text?};
  % n469 - Formula "\frac{a}{b}"
  \node[fcformula, scale=1.138, text=fc12161C]
        (n469) at (7.46,-2.8) {$\displaystyle \frac{a}{b}$};
  % n470 - Chemical formula "H2O"
  \node[fcformula, scale=1.138, text=fc12161C] (n470) at (7.4,-4.36) {\ce{H2O}};
  % n471 - Resistor "R"
  \draw[fccircuit, color=fc1F2933, line width=0.683pt, /tikz/circuitikz/bipoles/length=1.4cm]
        (6.8,-5.8) to[R, l={$R$}] (8.8,-5.8);
  % n472 - Molecule "3 atoms"
  \node[fcformula] (n472) at (11.2949,-2.876) {%
    \scalebox{0.797}{\color{fc12161C}\setchemfig{atom sep=0.853cm, bond style={draw=fc1F2933, line width=0.857pt}}\chemfig{{O}(-[:210]{H})-[:330]{H}}}};
  % ---- Connectors ----
  % e468 - n466 -> n467 (orthogonal)
  \draw[fcedge, -{Stealth[length=5.12pt, width=3.84pt]}, draw=fc1F2933, line width=0.91pt]
    (3.6,-3.4) -- (3.6,-3.84) -- (3.6,-4.1) -- (3.5,-4.1) -- (3.5,-4.36) -- (3.5,-4.8);
\end{tikzpicture}
\end{document}